\section{Hash Layers}

RCOUNT uses separate hash roles so that verifier transcripts can name the layer
that failed.

\begin{center}
\small
\begin{tabularx}{\textwidth}{lXX}
\toprule
Layer & Object & Purpose \\
\midrule
Source hash & Raw public input file & Detects byte changes in the evidence an adapter read. \\
Record hash & Normalized record & Supports stable hashes over canonicalized rows. \\
Event hash & Status or lineage event & Binds typed public events to canonical records. \\
Package hash & RCOUNT package content & Names the normalized claim set as a whole. \\
Proof hash & Public proof object & Names inclusion-proof material without exposing choices. \\
\bottomrule
\end{tabularx}
\end{center}

The current crates use domain-separated SHA-256 prefixes for these roles. The
domain separation matters because a hash of a source file should not be confused
with a hash of a package, event, or proof. Each role answers a different
question.

\begin{center}
\small
\begin{tabularx}{\textwidth}{lX}
\toprule
Prefix & Hash role \\
\midrule
\texttt{RCOUNT\_SOURCE\_V1} & Raw source file bytes. \\
\texttt{RCOUNT\_RECORD\_V1} & Normalized canonical record. \\
\texttt{RCOUNT\_FILE\_V1} & Canonical normalized file projection. \\
\texttt{RCOUNT\_PACKAGE\_V1} & Package content hash. \\
\texttt{RCOUNT\_EVENT\_V1} & Status or canvass event. \\
\texttt{RCOUNT\_PROOF\_V1} & Public proof object. \\
\bottomrule
\end{tabularx}
\end{center}

The suffix is part of the contract. Future RCOUNT versions may migrate hash
algorithms or canonical projections, but the prefix and algorithm must be
recorded so older packages remain replayable.

The \texttt{summary-basic} package currently records a package content hash and
simple package counts:
\begin{quote}
\footnotesize
\begin{verbatim}
{
  "package_content_hash": "sha256:2e8876...",
  "contest_count": 1,
  "reporting_unit_count": 3,
  "batch_count": 0,
  "summary_count": 3
}
\end{verbatim}
\end{quote}

That digest names the normalized package. It does not replace the source index:
an auditor still needs to know which raw files were read and whether those raw
file hashes match.

The normalized package hash covers the package's contests, reporting units,
batches, lineage records, inclusion-proof records, summaries, and status events
after canonical JSON projection. Volatile metadata such as local file mtimes,
temporary paths, pretty-printing, and transcript generation time is not part of
that normalized package content.
